\documentclass[a4paper,12pt]{article}
\usepackage[spanish, es-nodecimaldot,mexico]{babel}

% *** FUENTE ***
\usepackage[T1]{fontenc}
\usepackage[lining,light]{InriaSans}
\renewcommand*\oldstylenums[1]{{\fontfamily{InriaSans-OsF}\selectfont #1}}
%% The font package uses mweights.sty which has som issues with the
%% \normalfont command. The following two lines fixes this issue.
\let\oldnormalfont\normalfont
\def\normalfont{\oldnormalfont\mdseries}

%*** GEOMETRÍA ***
\usepackage{geometry}
\geometry{left=25mm,right=25mm,top=35mm,bottom=30mm,headheight=35mm,showframe}

%*** FECHA ***
\usepackage[useregional]{datetime2}

% *** COLOR ***
\usepackage[table]{xcolor}
\definecolor{rojo}{RGB}{217, 62, 37}
\definecolor{gris}{RGB}{242, 242, 242}  % Corregido: 242,242,242

% *** FIGURAS ***
\usepackage{graphicx} %Para cargar figuras
\graphicspath{{./Figuras}} % inidican donde están las figuras
\usepackage{float} % Permite usar opción H
\usepackage{subcaption}  %Para poner subfiguras
\usepackage{wrapfig} % Para poner texto a lado de las figuras
\usepackage{overpic}
\usepackage{caption}

% *** ECUACIONES ***
\usepackage{mathtools}

% *** REFERENCIAS ***
\usepackage{hyperref}
\hypersetup{colorlinks=true, linkbordercolor=blue}

\begin{document}
hola

\begin{minipage} {0.45\linewidth}
    \includegraphics[width=1\linewidth]{example-image-a}
\end{minipage}
\begin{minipage} {0.45\linewidth}
    \includegraphics[width=1\linewidth]{example-image-b}
\end{minipage}

\begin{center}
\begin{minipage} {0.5\linewidth}
    \includegraphics[width=1\linewidth]{example-image-a}
    \captionof{figure}{Leyenda de la Figura A}
\end{minipage} \hfil
\begin{minipage} {0.3\linewidth}
    At a minimum, you’ll need a TeX distribution, a good text editor and a DVI or PDF viewer.
More specifically, the basic requirement is to have a TeX compiler (which is used to generate
output files from source), fonts, and the LaTeX macro set.
\end{minipage}
\end{center}

\section{Ecuaciones}
Vamos a poner un poco de texto áca  para que se note que sabemos lo que decimos.
Ecuación. 
\begin{center}
    $ E = m \cdot c^2 $. Tenemos $ f_{yk} $ o sino $ x_{a}^{2}$. También tenemos 
    $ \frac{3}{2} $ o si no $ \dfrac{3}{4} $ o la $ \sqrt{2} $ o tambien 
    $ \sqrt[3]{9} $
\end{center} 

$$ E = m \cdot c^2 $$

\begin{equation}
    E = m \cdot c^2
     \label{eq:energia}
\end{equation}

Como se muestra en la ecuación \ref{eq:energia}. La \autoref{eq:energia}
o si nola ecuación \eqref{eq:energia}
La \

\begin{equation*}
    E = m \cdot c^2
\end{equation*}

La suma es $ a+b $ y la resta es $ a-b $ Multiplicación es $ a \cdot b $
o sino $ a \times  b\  \text{hola}$
Si estoy con ángulos puedo tener $ \cos(30) $ o sino $ \sin(30) $.

Otra opción es una ecuación su numerar pero dentro del entorno de
ecauaciones. 

Una integral"
$$ \int_0^L x^2, \quad 10<15 \Rightarrow 15 \geq  10$$


\begin{align}
    f(x) & = x^2+x^3. \\
         & = x^2+x^2 \cdot x
\end{align}

\begin{align*}
    f(x) & = x^2+x^3. \\
         & = x^2+x^2 \cdot x
\end{align*}


\begin{align*}
    A &=
    \begin{pmatrix}
        1 & 2 & 3 \\
        4 & 5 & 6 \\
        7 & 8 & 9
    \end{pmatrix}, \\[6pt]
    B &=
    \begin{bmatrix}
        1 & 2 & 3 \\
        4 & 5 & 6 \\
        7 & 8 & 9
    \end{bmatrix}, \\[6pt]
    C &=
    \begin{matrix}
        1 & 2 & 3 \\
        4 & 5 & 6 \\
        7 & 8 & 9
    \end{matrix}, \\[6pt]
    D &=
    \begin{Bmatrix}
        1 & 2 & 3 \\
        4 & 5 & 6 \\
        7 & 8 & 9
    \end{Bmatrix}, \\[6pt]
    E &=
    \begin{vmatrix}
        a & b_2 & 3 \\
        4 & 5 & 6 \\
        7 & 8 & 9
    \end{vmatrix}, \\[6pt]
    F &=
    \begin{Vmatrix}
        1 & 2 & 3 \\
        4 & 5 & 6 \\
        7 & 8 & 9
    \end{Vmatrix},
\end{align*} 

\begin{equation}
\sum _{x=0}^{n}f{\left(x\right)}^{2}+1=0
\end{equation}

\end{document}